\documentclass[sigconf]{acmart}

\setcopyright{none}
\settopmatter{printacmref=false, printfolios=true}
\renewcommand\footnotetextcopyrightpermission[1]{}
\pagestyle{plain}

\usepackage{booktabs}
\usepackage{amsmath}
\usepackage{balance}

\acmConference[ACM SIGSPATIAL '26]{The 34th ACM International Conference on Advances in Geographic Information Systems}{November 3--6, 2026}{Riverside, CA, USA}
\acmYear{2026}
\copyrightyear{2026}

\newcommand{\conf}{\mathcal{C}_{\tau}}
\newcommand{\dview}{\Delta_{\mathrm{view}}}

\title{Conflict-Aware Cross-View Disaster Triage Across Spatial Observation Regimes [experiments]}

\author{Yifan Yang}
\affiliation{%
  \department{Department of Geography}
  \institution{Texas A\&M University}
  \city{College Station}
  \state{Texas}
  \country{USA}
}

\author{Lei Zou}
\affiliation{%
  \department{Department of Geography}
  \institution{Texas A\&M University}
  \city{College Station}
  \state{Texas}
  \country{USA}
}

\begin{abstract}
Cross-view disaster assessment combines ground-level and overhead imagery to
infer building damage. Aggregate F1 scores, however, hide the cases where
fusion matters most: examples for which the single-view predictors disagree.
This short paper introduces a Conflict-Aware Evaluation protocol (CAE) for
paired-view disaster triage. CAE reports overall F1, conflict-subset F1, and
the conflict gain of cross-view fusion over the stronger single-view baseline.
We apply CAE to two paired disaster settings: an Eaton/Altadena wildfire dataset
with property-centered ground imagery and the Hurricane Ian CVDisaster benchmark
with panoramic ground imagery. The results show that cross-view fusion is most
informative on disagreement cases, but the dominant single-view arbitrator
depends on the observation regime and the split definition. On the wildfire
sensitivity split, cross-view fusion improves conflict F1 by 0.0388 over the
best single-view baseline. On a leakage-controlled hurricane endpoint split,
cross-view fusion no longer dominates street-only inference on conflict cases,
which provides an important boundary condition. We further report two pilot
extensions: conflict density as an unsupervised tile-level wildfire damage
proxy and a conflict-weighted loss that improves wildfire conflict F1. The
evidence supports a conservative conclusion: paired-view disaster models should
be evaluated as conflict resolvers, not only as average-case classifiers.
\end{abstract}

\keywords{disaster assessment, cross-view learning, multimodal fusion, GeoAI, conflict-aware evaluation}

\begin{document}

\maketitle

\section{Introduction}

Post-disaster building assessment increasingly uses several visual viewpoints.
Overhead imagery gives broad spatial coverage, while ground imagery can expose
facade condition, debris, and local structural cues. Cross-view models seek to
combine these signals, but prior evaluation often reduces their value to one
aggregate test metric. That view is too coarse for operational disaster triage.
Many buildings are easy in both views. The harder and more consequential cases
are those where the two views provide competing evidence.

This paper studies cross-view fusion as a conflict-resolution problem. We focus
on paired examples that contain a ground image, an overhead patch, and a binary
damage label. Instead of asking only whether paired-view fusion improves
average F1, we ask when fusion helps after the two single-view models disagree.
That framing matters for geographic AI because observation conditions are not
interchangeable. A close property inspection photograph and a panoramic street
scene both count as ground imagery, but they encode different relationships
between observer, target building, and surrounding landscape.

The study sits between several lines of disaster-vision work. Satellite damage
benchmarks such as xBD have made overhead-only building assessment a standard
task \cite{gupta2019creating}, and later fusion models improve satellite-based
damage assessment from multi-scale visual evidence \cite{xiao2023cross}. Ground
and street-view studies show that close-range imagery can capture damage cues
that overhead imagery misses \cite{zhai2020damage,xue2024post,yang2025hyperlocal}.
CVDisaster is the closest paired-view benchmark for Hurricane Ian
\cite{li2025cross}. CAE adds a different evaluation axis by separating average
performance from behavior on single-view disagreement cases.

We contribute a compact evaluation study for ACM SIGSPATIAL short-paper review.
First, we define Conflict-Aware Evaluation (CAE), a three-metric protocol that
reports full-test performance and conflict-subset behavior together. Second, we
apply CAE to wildfire and hurricane paired-view triage, while treating the
hurricane grouped-clean split as the conservative endpoint control. Third, we
show how the same disagreement signal can support two extensions: wildfire
tile-level conflict-density mapping and a conflict-weighted training loss. The
paper deliberately avoids a causal claim that disaster type alone determines
which view should dominate. The evidence is better read as a regime-sensitive
experiment: cross-view fusion helps most when single-view evidence conflicts,
and the source of that help depends on how the ground view observes the target.

\section{Data and Protocol}

\subsection{Paired-View Triage Settings}

Let
$\mathcal{D}=\{(x_i^{(s)},x_i^{(r)},y_i)\}_{i=1}^{N}$ be a paired disaster
triage dataset, where $x_i^{(s)}$ is a street or ground image,
$x_i^{(r)}$ is an overhead patch, and $y_i\in\{0,1\}$ is the binary damage
label. We compare three modes: \emph{street-only}, \emph{remote-only}, and
\emph{cross-view}. The cross-view model uses late fusion of ResNet18 features
\cite{he2016deep}, including each embedding, their absolute difference, and
their elementwise product. This simple fusion design keeps the evaluation
focused on conflict behavior rather than architectural novelty.

The wildfire setting uses an Eaton/Altadena paired dataset assembled for this
project. Its ground imagery is property-centered inspection imagery paired with
overhead patches. This observation regime often places the target building near
the visual center of the ground image. The hurricane setting uses the Hurricane
Ian CVDisaster benchmark \cite{li2025cross}, which pairs panoramic ground
imagery with very high resolution overhead imagery. This observation regime
often includes wider street context, adjacent buildings, and weaker target
alignment. The comparison therefore mixes disaster type, geography, source, and
collection protocol; those factors limit causal interpretation and motivate the
conservative language used below.

\subsection{Conflict-Aware Evaluation}

For each example, let $p_i^{(s)}$ and $p_i^{(r)}$ denote the positive-class
probabilities predicted by the street-only and remote-only models. CAE defines
the conflict subset
\begin{equation}
\conf=\{i: |p_i^{(s)}-p_i^{(r)}|>\tau\},
\end{equation}
with $\tau=0.1$ in the main analysis. For each model $m$, CAE reports standard
overall F1 and conflict F1 on $\conf$. It then computes the cross-view conflict
gain
\begin{equation}
\dview =
\mathrm{F1}(\mathrm{crossview};\conf)
-
\max\{\mathrm{F1}(\mathrm{street};\conf),
\mathrm{F1}(\mathrm{remote};\conf)\}.
\end{equation}
This protocol exposes whether paired-view fusion resolves disagreement cases
rather than only improving easy average cases.

We also decompose conflicts by correction direction. $S{\rightarrow}O$ denotes
cases where the street-only model is correct and the remote-only model is
wrong; $O{\rightarrow}S$ denotes the reverse. These fractions identify the
stronger single-view arbitrator inside the conflict subset. Because a prior
hurricane endpoint split had object-level leakage, we report the grouped-clean
endpoint split as a required control rather than as an optional robustness run.

\section{Results}

\subsection{CAE Summary}

Table~\ref{tab:cae} reports the core CAE results. The wildfire sensitivity
split has the clearest positive conflict gain: cross-view fusion reaches 0.6618
conflict F1, which is 0.0388 above the better single-view baseline. The original
hurricane endpoint split also shows a large gain, but its small conflict subset
($n=32$) and leakage history make it diagnostic rather than primary evidence.
The broader hurricane moderate-plus-severe split keeps a smaller positive gain.
The grouped-clean endpoint split provides the most important boundary condition:
cross-view fusion remains close to street-only overall, but it no longer
dominates street-only on the conflict subset.

\begin{table}[t]
  \caption{Conflict-Aware Evaluation summary. Conflict F1 uses $\conf$ with
  $\tau=0.1$. WF is wildfire; H is hurricane. Gain is reported only for
  cross-view rows.}
  \label{tab:cae}
  \centering
  \scriptsize
  \setlength{\tabcolsep}{2.6pt}
  \begin{tabular}{llrrrr}
    \toprule
    Setting & M & All & C-F1 & Gain & $n_c$ \\
    \midrule
    WF sens. & S & 0.9392 & 0.6230 & -- & 435 \\
    WF sens. & R & 0.9238 & 0.5237 & -- & 435 \\
    WF sens. & CV & 0.9473 & 0.6618 & +0.0388 & 435 \\
    H end. & S & 0.9055 & 0.6842 & -- & 32 \\
    H end. & R & 0.8969 & 0.5806 & -- & 32 \\
    H end. & CV & 0.9192 & 0.7647 & +0.0805 & 32 \\
    H mod.+sev. & S & 0.8738 & 0.7324 & -- & 128 \\
    H mod.+sev. & R & 0.8873 & 0.7742 & -- & 128 \\
    H mod.+sev. & CV & 0.8915 & 0.7867 & +0.0125 & 128 \\
    H clean & S & 0.9016 & 0.8101 & -- & 80 \\
    H clean & R & 0.8385 & 0.6526 & -- & 80 \\
    H clean & CV & 0.9008 & 0.7949 & -0.0153 & 80 \\
    \bottomrule
  \end{tabular}
\end{table}

The directional decomposition sharpens this interpretation. On wildfire,
$S{\rightarrow}O$ conflicts are more common than $O{\rightarrow}S$ conflicts
(0.2713 versus 0.1885), so the ground image is the stronger single-view
arbitrator. On the hurricane endpoint split, the two directions are near tied
(0.3750 versus 0.3438). On the broader hurricane split, the ordering reverses
(0.1094 versus 0.1328), making the overhead view the stronger single-view
arbitrator. The stable claim is therefore not that one disaster always favors
one modality. It is that conflict behavior changes with the observation regime
and split definition.

\subsection{Spatial and Training Extensions}

CAE also produces useful secondary signals. For the wildfire dataset, we
aggregate $|p_i^{(s)}-p_i^{(r)}|$ by overhead tile and compare the resulting
conflict density with observed tile damage rate. Across 27 tiles, Spearman
$r=0.5123$ and $p=0.0063$. This suggests that cross-view disagreement can act
as an unsupervised spatial damage proxy in the property-centered wildfire
setting. We do not yet claim the same result for hurricane because the benchmark
does not provide equivalent tile metadata in this repository.

We also test ConflictFocalLoss, a conflict-weighted binary loss related in
spirit to focal reweighting \cite{lin2017focal}. For each cross-view example,
the loss multiplies binary cross entropy by $(1+c_i)$, where $c_i$ is a
feature-space disagreement score based on the cosine distance between the two
view embeddings. On the wildfire sensitivity split, moderate conflict
weighting improves overall F1 from 0.9473 to 0.9520 and conflict F1 from
0.6618 to 0.7208. This result is encouraging, but it is still a single-setting
pilot rather than a fully validated method contribution.

\begin{table}[t]
  \caption{Pilot method and spatial extensions.}
  \label{tab:extensions}
  \centering
  \small
  \setlength{\tabcolsep}{4pt}
  \begin{tabular}{lrr}
    \toprule
    Extension & Primary statistic & Scope \\
    \midrule
    Conflict density & $r=0.5123$, $p=0.0063$ & wildfire tiles \\
    ConflictFocalLoss & conflict F1 $=0.7208$ & wildfire split \\
    Adaptive cascade & F1 $=0.9498$, 9.5\% routed & wildfire split \\
    \bottomrule
  \end{tabular}
\end{table}

\section{Discussion and Limitations}

The experiment supports CAE as a useful evaluation protocol for paired-view
disaster triage. Aggregate F1 alone would make the models look similar in some
settings. CAE reveals how they behave on disagreement cases and shows that
cross-view value can shrink, grow, or even become negative relative to the
stronger single-view model after leakage control.

Several limits are central. First, the wildfire and hurricane settings differ
by more than observation regime. They also differ by hazard, geography, image
source, label construction, and metadata availability. The results therefore
support a regime-sensitive interpretation, not a causal proof. Second, the
hurricane endpoint conflict subset is small, so its large positive gain should
be treated as tentative. Third, the private wildfire data package must include
complete counts, class balance, annotation protocol, inter-annotator agreement
if available, and release or access instructions before final submission.
Fourth, ConflictFocalLoss and the tile-level spatial proxy are validated only
on the wildfire setting. They motivate follow-up experiments but should not be
presented as fully general.

These limits also clarify the paper's contribution. CAE does not claim that
cross-view fusion always wins. It gives reviewers and practitioners a clearer
way to see when paired views resolve evidence conflict and when a single view
remains sufficient.

\section{Conclusion}

This short paper reframes paired-view disaster triage as conflict-aware
evaluation. Across wildfire and hurricane paired-view settings, cross-view
fusion is most informative on disagreement cases, but the stronger single-view
arbitrator changes with observation regime, split definition, and leakage
control. The resulting lesson is practical: disaster benchmarks should report
conflict-subset behavior alongside aggregate scores. Doing so makes multimodal
fusion more falsifiable, more spatially interpretable, and more useful for
deployment decisions.

\begin{acks}
The authors used AI-based writing assistance to support grammar checking,
draft compression, and editorial revision. These tools were not used as sources
of scientific evidence or to autonomously generate the study design, datasets,
experimental results, figures, tables, interpretations, or conclusions. The
authors reviewed and approved all manuscript content and take full
responsibility for the work.
\end{acks}

\balance
\bibliographystyle{ACM-Reference-Format}
\bibliography{references}

\end{document}
